% ===================================================================
%  TAFEL Nr. 11 — Die Stadt und ihre Gebäude. --- Der Brand.
%  Traduction 2026 d'apres texto/io, controlee sur texto/fr et
%  texto/en. Consigne : voir texto/de/10-tafel-01.tex.
%
%  QUATRE INSTITUTIONS SUR UNE MEME PAGE, ET CHACUNE PARLE LA LANGUE
%  DU PAYS QUI L'A INVENTEE.
%
%  1. LA CAISSE D'EPARGNE, ALINEA 6 — Sparkasse. Fondee a
%     Hambourg en 1778, nommee en allemand, compose transparent :
%     la caisse ou l'on epargne. Le tableau 11 romanche avait deja
%     bute sur cet objet et trouve un compose lui aussi.
%
%  2. LES POMPIERS, SECONDE SCENE — Feuerwehr, la
%     defense-contre-le-feu. Les corps de volontaires sont nes a
%     Durlach en 1846, et tout leur materiel porte un compose
%     allemand fait a cette date : Feuerwehrmann, Feuerspritze (67),
%     Handspritze (68), Schlauch (69), Strahlrohr (70), Wasserstrahl
%     (71), Feuerleiter (72), Feuermelder. C'est la cinquieme campagne
%     de nomination de la colonne, apres Jahn, le Sprachverein, la
%     marine — et c'est la plus complete des quatre, parce que
%     l'institution etait NEUVE : il n'y avait rien a remplacer. Le
%     tableau 11 estonien avait vu son « tuletõrjujad » faire
%     exactement le meme compose, la aussi pour un corps fonde de
%     toutes pieces.
%
%  3. LE THEATRE ET L'OPERA, ALINEAS 5 A 7 — et tout y est italien :
%     Theater, Oper, Foyer, Loge, Parkett, Orchester, Kulisse,
%     Souffleur, Tenor, Kostüm, Portikus. L'opera est ne a Florence
%     vers 1600 et il a voyage avec son lexique entier, comme le
%     glacier du tableau 9.
%
%  4. LA BANQUE, ALINEA 10 — Bank. Le mot est italien, de
%     « banca », LE BANC du changeur sur le marche. Or « Bank » est
%     aussi le banc allemand, le meme meuble, le meme mot germanique
%     que l'italien avait emprunte au lombard mille ans plus tot. LE
%     MOT EST PARTI EN ITALIE COMME MEUBLE ET REVENU EN ALLEMAGNE
%     COMME METIER. C'est le seul cas de la colonne ou un emprunt
%     soit un retour.
%
%  D'OU LA REGLE, QUI VAUT POUR LES ONZE TABLEAUX PRECEDENTS :
%
%     LA LANGUE D'UNE INSTITUTION EST CELLE DU PAYS OU ON L'A BATIE
%     LA PREMIERE FOIS.
%
%  Elle explique les corporations allemandes du tableau 5 (Allemagne,
%  XIIIe siecle), la vigne latine du tableau 8 (Rome), les grades
%  latins du tableau 10 (Rome encore), le gymnase de Jahn et la
%  Feuerwehr de Durlach. Elle explique aussi ce que l'allemand
%  N'A PAS traduit : Präfektur, Kommissar, Journalist, Ambulanz,
%  Boulevard — la prefecture, le commissaire, le journal, l'ambulance
%  et le boulevard sont des institutions francaises, et la ville de
%  cette planche est une ville francaise.
%
%  ET UN BLOC QUI PIEGE DEUX COLONNES SUR TRENTE, POUR LA MEME
%  RAISON. Au bloc c2-02-1, l'ido range les soldats (57), les
%  gendarmes (63), puis les curieux (56) : le comparatif porte sur le
%  VERBE, et l'objet vient apres. L'allemand veut son objet avant
%  « besser … als », et le romanche voulait le sien avant « meglier …
%  che » : les deux colonnes ont fait la meme inversion au meme
%  endroit, sans aucun rapport de langue. LE PIEGE N'EST PAS DANS LA
%  LANGUE D'ARRIVEE, IL EST DANS LA PHRASE DE DEPART. Meme remede
%  dans les deux : faire porter le comparatif sur le verbe et rejeter
%  l'objet dans une seconde proposition.
%
%  UNE NOTE DE COMPOSITION. L'alinea 9 de la premiere scene porte, en
%  ido comme en francais, un gras errant : « quale », « La
%  vetera veturi », « basa prezo ». Le rendu ote deja le gras des
%  mots-outils ; on ecrit ici l'allemand comme il se dit, et le gras
%  tombe ou il doit.
% ===================================================================

%%K t11-apar-1 apar
\VUsaut{2.0mm}
\VUcentre{13.2pt}{90}{DRITTE REIHE}
\VUsaut{3.0mm}
\VUfilet{16mm}
\VUsaut{5.6mm}
\VUtitre{15.0pt}{60}{TAFEL Nr. 11}
\VUsaut{1.4mm}
\VUpk{11.4pt}{40}{Die Stadt und ihre Gebäude. --- Der Brand.}
\VUsaut{2.2mm}

%%K t11-tit-1 sub
\VUpk{10.2pt}{40}{Die Vorstadt.}
\VUsaut{3.0mm}

%%K t11-c1-01-1 p
1. --- Ich wohne in einer großen Stadt hundert Kilometer vom Meer, die trotzdem ein
erstklassiger \VUgras{Hafen} \textsuperscript{(1)} ist. Der Fluss, der an ihr
entlangfließt, ist hier vierhundert Meter breit und bildet eine sehr schöne Reede.

%%K t11-c1-02-1 p
2. --- Der ganze Teil der Stadt an den \VUgras{Kais} \textsuperscript{(2)} sieht ganz
nach Seefahrt und Industrie aus und bildet eine \VUgras{Vorstadt}
\textsuperscript{(3)}, in der nur Seeleute und Arbeiter wohnen. Diese arbeiten in den
\VUgras{Fabriken} \textsuperscript{(4)} und im \VUgras{Gaswerk}
\textsuperscript{(5)}, dessen hoher \VUgras{Schornstein} \textsuperscript{(6)} und
dessen \VUgras{Gasbehälter} von weitem zu sehen sind.

%%K t11-c1-03-1 p
3. --- Im Mittelalter war die Stadt befestigt, und man findet noch Reste ihrer
Wälle, vor allem das \VUgras{Tor} \textsuperscript{(8)} der alten \VUgras{Burg}
\textsuperscript{(7)} mit ihren beiden \VUgras{Türmen} \textsuperscript{(9)}, die
heute als \VUgras{Gefängnis} dienen.

%%K t11-c1-04-1 p
4. --- In diesem ganzen \VUgras{Viertel}, das sehr alt aussieht, ist nur ein Gebäude
verhältnismäßig neu: das \VUgras{Zollamt} \textsuperscript{(10)}. Es steht zwischen
dem Kai und der kleinen \VUgras{Gasse} \textsuperscript{(11)}, die zum alten
\VUgras{Rathaus} \textsuperscript{(12)} führt. Dieses letzte Gebäude erkennt man
leicht an dem riesigen \VUgras{Belfried} \textsuperscript{(13)}, der über der
Altstadt aufragt.

%%K t11-c1-05-1 p
5. --- Auf der anderen Seite der Gasse \VUgras{steht} ein Gebäude im gleichen Stil
wie das Zollamt: es ist die \VUgras{Börse} \textsuperscript{(14)}. Eine ihrer Fronten
geht auf einen kleinen Platz, an dem die \VUgras{Präfektur} \textsuperscript{(15)}
liegt. Unsere Stadt ist keine Hauptstadt, aber doch ein wichtiger Bezirksvorort.

%%K t11-tit-2 sub
\VUsaut{5.0mm}
\VUpk{10.2pt}{40}{Die Oberstadt.}
\VUsaut{3.4mm}

%%K t11-c1-06-1 p
6. --- Die Garnison, eine ziemlich große, ist in weiten und sehr gesunden
\VUgras{Kasernen} \textsuperscript{(16)} untergebracht; sie reichen bis an das Gitter
des \VUgras{Stadtparks} \textsuperscript{(17)}. Der \VUgras{Boulevard}
\textsuperscript{(19)}, der zu diesem Park führt, läuft hinter der
\VUgras{Sparkasse} \textsuperscript{(18)} entlang und mündet auf den Platz vor dem
\VUgras{Dom} \textsuperscript{(20)}.

%%K t11-c1-07-1 p
7. --- All diese Gebäude steigen in Stufen auf einem kleinen Hügel an, auf dem auch
unser großes \VUgras{Gymnasium} \textsuperscript{(21)} steht.

%%K t11-c1-07-2 p
Geht man eine leicht abschüssige Straße hinunter, die auf die Hauptstraße trifft, so
kommt man zum \VUgras{Gerichtsgebäude} \textsuperscript{(22)}, vor dem eine
angenehme \VUgras{Grünanlage} \textsuperscript{(26)} liegt. Dieser \VUgras{Platz} ist
nicht sehr groß, aber mitten in der Stadt bildet er eine Oase aus Grün und kühler
Luft. Die Stadtleute nutzen ihn deshalb viel und setzen sich gern unter die Markise
des \VUgras{Cafés} \textsuperscript{(25)}, um der Militärkapelle zuzuhören, die
donnerstags und sonntags im \VUgras{Musikpavillon} \textsuperscript{(27)} zwischen
den Rasenflächen und den Blumenbeeten spielt.

%%K t11-c1-08-1 p
8. --- Auf einer dieser Rasenflächen steht eine bronzene \VUgras{Statue}
\textsuperscript{(28)}, ein Denkmal zur Erinnerung an einen unserer Mitbürger, der
sich um seine Geburtsstadt große Verdienste erworben hat.

%%K t11-tit-3 sub
\VUsaut{4.0mm}
\VUpk{10.2pt}{40}{Unterwegs in der Stadt.}
\VUsaut{3.4mm}

%%K t11-c1-09-1 p
9. --- Unsere Stadt hat noch kein elektrisches Licht; sie hat nur
\VUgras{Gaslaternen} \textsuperscript{(29)}. Aber die \VUgras{örtliche Presse}
setzt sich dafür ein, dass diese Beleuchtung verbessert wird, \VUgras{so wie} man
den \VUgras{Antrieb} der Straßenbahn schon verbessert hat. Die alten
\VUgras{Pferdewagen} sind durch praktische \VUgras{elektrische Straßenbahnen}
\textsuperscript{(31)} ersetzt worden, die die Fahrgäste schnell von einem
\VUgras{Ende der Stadt zum anderen} bringen, zu einem sehr \VUgras{niedrigen}
Fahrpreis.

%%K t11-c1-10-1 p
10. --- Neben dem Gerichtsgebäude liegt die \VUgras{Bank} \textsuperscript{(32)}, wo
Handelswechsel diskontiert werden. Die \VUgras{Bankboten} \textsuperscript{(33)} gehen
zu den Leuten nach Hause, um das Fällige einzuziehen.

%%K t11-tit-4 sub
\VUsaut{4.0mm}
\VUpk{10.2pt}{40}{Kunst und Wissenschaft.}
\VUsaut{3.4mm}

%%K t11-c1-11-1 p
11. --- Das schöne Gebäude gleich daneben ist das \VUgras{Museum}. Es beherbergt
zugleich die \VUgras{Stadtbibliothek} \textsuperscript{(34)}, schöne
\VUgras{naturkundliche Sammlungen} \textsuperscript{(35)} (ein Naturkundemuseum) und
eine anmutige \VUgras{Gemäldegalerie} \textsuperscript{(36)}, die eine prächtige
ständige \VUgras{Ausstellung} bildet.

%%K t11-c1-11-2 p
Es ist zweimal in der Woche für alle offen, aber Besucher können es an jedem Tag
gegen ein kleines Trinkgeld an den \VUgras{Aufseher} \textsuperscript{(37)} ansehen.
\VUgras{Maler} \textsuperscript{(38)} kommen her, um zu arbeiten. Man sieht sie
\VUgras{vor} ihren Staffeleien stehen, die \VUgras{Palette} in der Hand, und die
berühmten Bilder kopieren.

%%K t11-c1-12-1 p
12. --- Auf den Höhen hinter dem Museum erkennt man gerade noch die \VUgras{Kuppel}
\textsuperscript{(40)} des \VUgras{Krankenhauses} \textsuperscript{(39)} und die
Gebäude des \VUgras{Lehrerseminars} \textsuperscript{(41)}, in dem die Lehrer unserer
Volksschulen ausgebildet worden sind. Am Theaterplatz gibt es ein \VUgras{Postamt}
\textsuperscript{(43)}, das in einem \VUgras{Privathaus} \textsuperscript{(42)}
untergebracht ist.

%%K t11-tit-5 sub
\VUsaut{5.0mm}
\VUpk{11.4pt}{40}{Der Brand.}
\VUsaut{3.0mm}

%%K t11-tit-6 sub
\VUpk{10.2pt}{40}{Feuerruf.}
\VUsaut{3.4mm}

%%K t11-c2-01-1 p
1. --- Eine beunruhigende Nachricht läuft durch die Stadt: im Theater ist Feuer
ausgebrochen! Als ich den \VUgras{Platz} \textsuperscript{(96)} erreiche, hat sich die
Menge schon auf dem \VUgras{Bürgersteig} \textsuperscript{(46)} und auf der
\VUgras{Fahrbahn} \textsuperscript{(47)} versammelt. Die \VUgras{Postbeamten}
\textsuperscript{(44)} und die \VUgras{Briefträger} \textsuperscript{(45)} kommen aus
dem Postamt. Ein \VUgras{Straßenbahnführer} \textsuperscript{(48)} hat seinen Wagen
anhalten müssen, um einen städtischen \VUgras{Krankenwagen} \textsuperscript{(50)}
durchzulassen, der mit einem \VUgras{Chirurgen} \textsuperscript{(52)} und einer
\VUgras{Krankenschwester} \textsuperscript{(51)} kommt, um sich um einen
\VUgras{Verletzten} \textsuperscript{(53)} zu kümmern, den man auf einer
\VUgras{Trage} \textsuperscript{(54)} an eine Stelle nahe dem \VUgras{Zeitungskiosk}
\textsuperscript{(55)} gebracht hat.

%%K t11-c2-02-1 p
2. --- Eine Kompanie Infanterie auf dem Rückweg vom Exerzieren hat angehalten, um
den berittenen \VUgras{Stadtgardisten} \textsuperscript{(61)} beim Ordnungshalten zu
helfen. Die \VUgras{Reiter}, deren Pferde sich aufbäumen, machen es besser als die
\VUgras{Soldaten} \textsuperscript{(57)} oder die \VUgras{Gendarmen}
\textsuperscript{(63)}: sie drängen die \VUgras{Zuschauer} \textsuperscript{(56)}
zurück. Die
Gendarmen haben ohnehin alle Hände voll zu tun. Einer von ihnen sagt den
\VUgras{Polizisten} \textsuperscript{(64)}, wohin sie einen zweiten
\VUgras{Verletzten} \textsuperscript{(65)} bringen sollen, einen tapferen
Feuerwehrmann, der von einer Leiter gestürzt ist.

%%K t11-c2-03-1 p
3. --- Der \VUgras{Offizier} \textsuperscript{(66)} zeigt genau, wo die
\VUgras{Feuerspritze} \textsuperscript{(67)}, die eben ankommt, stehen soll. Solange
er auf diese starke Maschine wartet, hat er eine \VUgras{Handspritze}
\textsuperscript{(68)} aufstellen lassen, deren langer \VUgras{Schlauch}
\textsuperscript{(69)} Ströme von Wasser auf das Feuer gießt. Sechs Feuerwehrleute
bedienen sie, während ihr Kamerad das \VUgras{Strahlrohr} \textsuperscript{(70)} hält
und den \VUgras{Wasserstrahl} \textsuperscript{(71)} mitten in den Brand richtet.

%%K t11-c2-04-1 p
4. --- Auf der anderen Seite des Theaters ist die große \VUgras{Leiter}
\textsuperscript{(72)} aufgerichtet worden. Sie lässt die Feuerwehrleute nahe an die
Flammen heran, die zwischen Schwaden schwarzen \VUgras{Rauchs}
\textsuperscript{(75)} emporschlagen.

%%K t11-tit-7 sub
\VUsaut{5.0mm}
\VUpk{10.2pt}{40}{Tapfere Taten.}
\VUsaut{3.4mm}

%%K t11-c2-05-1 p
5. --- Das \VUgras{Feuer} \textsuperscript{(74)} ist im Foyer des \VUgras{Theaters}
\textsuperscript{(73)} ausgebrochen, während einer Probe der \VUgras{Oper}, deren
erste \VUgras{Aufführung} morgen sein sollte. Zum Glück waren nur wenige
\VUgras{Zuschauer} \textsuperscript{(76)} im Saal. Die armen Leute haben sich auf den
\VUgras{Balkon} \textsuperscript{(77)} geflüchtet, wo tapfere
\VUgras{Feuerwehrleute} \textsuperscript{(86)} ihnen mit einer Leiter und mit Seilen
zu Hilfe kommen.

%%K t11-c2-06-1 p
6. --- Unter dem \VUgras{Portikus} \textsuperscript{(78)}, wo der \VUgras{Anschlag}
\textsuperscript{(79)} die Aufführung ankündigt, sieht man die \VUgras{Musiker}
\textsuperscript{(81)} des \VUgras{Orchesters} \textsuperscript{(80)} fliehen und ihre
Instrumente forttragen: \VUgras{Geige} \textsuperscript{(81)}, \VUgras{Flöte}
\textsuperscript{(82)}, \VUgras{Cello} \textsuperscript{(83)}, \VUgras{Posaune}
\textsuperscript{(84)}, \VUgras{Trommel} \textsuperscript{(85)} und so weiter. Man
versucht, einen Teil des \VUgras{Mobiliars} \textsuperscript{(87)} zu retten.

%%K t11-c2-07-1 p
7. --- Die \VUgras{Schauspieler} \textsuperscript{(89)} und
\VUgras{Schauspielerinnen} \textsuperscript{(88)}, \VUgras{Sänger} und Sängerinnen,
haben in ihren \VUgras{Bühnenkostümen} \textsuperscript{(90)} hinauslaufen müssen. Der
Tenor erklärt einem \VUgras{Journalisten} \textsuperscript{(92)}, wie es gekommen ist.
\VUgras{Der Reporter} ist vom ersten Augenblick an herbeigelaufen, zusammen mit den
Behörden: dem \VUgras{Präfekten} \textsuperscript{(91)}, dem
\VUgras{Bürgermeister} \textsuperscript{(93)}, dem \VUgras{Polizeikommissar}
\textsuperscript{(94)} und einem \VUgras{Untersuchungsrichter}
\textsuperscript{(95)}, der eine Untersuchung führen wird, um festzustellen, wer die
Verantwortung trägt.
\VUsaut{6.0mm}
\VUfilet{22mm}
